\section{Superseded and Auxiliary Validation and Limit Studies}
\label{app:prior-validation-results}

This appendix preserves the development history of the analysis without assigning the
older plots the status of v3 result figures. In particular, the closure panels below
predate the selected year-dependent lower-length-scale factors, while the limit figures
predate the July 8, 2026 configuration or use plotting repairs that have not yet been
replaced by rerun likelihood outputs. They remain useful for method auditing, but they
must not be combined numerically with the main-text candidate. Each subsection records
the question asked by the study, the relevant configuration, and the conclusion that
was carried forward. Captions deliberately retain the historical confidence level and
toy convention, even where these differ from the v3 90\% CL analysis.

\subsection{Earlier functional-form closure summaries}

Figures~\ref{fig:closure-validation-2015-funcform}--%
\ref{fig:closure-validation-2021-funcform} are the former main-text four-panel
functional-form summaries. They motivated the expanded lower-bound study, but they do
not implement the selected three-dataset choices summarized in
Table~\ref{tab:results-v3-configuration}. The error-barred, matched-reference
lower-bound scans in
Figures~\ref{fig:v3-lslb-closure-2015}--\ref{fig:v3-lslb-closure-2021} supersede
them as the headline closure evidence. The functional-form seed and toy-mean overview
is now shown where those pseudoexperiments are introduced in
Section~\ref{sec:toys}; the plots below retain the earlier extraction
diagnostics for comparison with that history.

\begin{figure}[p]
\centering
\graphicorplaceholder{0.99\linewidth}{toy_generation_figs/closure_validation_suite_2015_fShiftSigPowTail.png}
\caption{Historical 2015 \texttt{fShiftSigPowTail} functional-form, full-refit closure
suite for the $2.25\,\sigma_m$ blind and matched-reference geometry. The rows were
selected from the \texttt{wide\_refmatched} subset of the blind-geometry comparison,
so the injected amplitude was defined using the same refit convention used to assess
the widened window. From left to right, the panels show the ensemble pull mean, pull
width, recovered amplitude $\hat A/A_{\mathrm{inj}}$, and residual significance
$\hat Z-Z_{\mathrm{inj}}$ versus mass; the error bars are standard errors of the
ensemble mean or width. The near-zero background-only pull and $\Delta Z$ supported
the absence of a mass-coherent spurious signal, while the small negative shift at the
largest injection illustrated limited signal absorption. This study is retained for
provenance; the 100-toy lower-bound comparison in
Figure~\ref{fig:v3-lslb-closure-2015} is the v3 selection evidence.}
\label{fig:closure-validation-2015-funcform}
\end{figure}

\begin{figure}[p]
\centering
\graphicorplaceholder{0.99\linewidth}{toy_generation_figs/closure_validation_suite_2016_fShiftSigPowTail.png}
\caption{Historical 2016 10\% \texttt{fShiftSigPowTail} functional-form, full-refit
closure suite from \texttt{2016\_gpr/closure\_final}. A $2.25\,\sigma_m$ blind and
training exclusion and matched-refit signal normalization were used at each mass.
The pull mean and width test bias and uncertainty calibration, the recovered-amplitude
panel tests linearity, and $\hat Z-Z_{\mathrm{inj}}$ tests significance calibration
after repeating the local GP refit and extraction on every toy. The isolated behavior
of the $Z_{\mathrm{inj}}=1$ amplitude ratio was interpreted as a small-signal,
finite-ensemble effect rather than a large bias because the higher injections closed
more closely. This plot predates, and by itself does not establish, the v3 choice
$k_{\min}=0.9$; that choice is supported by
Figure~\ref{fig:v3-lslb-closure-2016}.}
\label{fig:closure-validation-2016-funcform}
\end{figure}

\begin{figure}[p]
\centering
\graphicorplaceholder{0.99\linewidth}{toy_generation_figs/closure_validation_suite_2021_fSigPowExpQ.png}
\caption{Historical 2021 1\% \texttt{fSigPowExpQ} functional-form, full-refit closure
suite from the \texttt{2021\_gpr/closure\_1pct/injection\_flat} toys. The study used
the then-current 2021 functional family and $2.25\,\sigma_m$ blind and training
exclusion. Pull mean and width assess background-only calibration; recovered amplitude
and $\hat Z-Z_{\mathrm{inj}}$ assess linearity and significance calibration for
$Z_{\mathrm{inj}}=1,3,5$. Pull means were near zero and widths below unity indicated
that the fitted uncertainty was larger than the ensemble scatter in this particular
toy model. The displayed support and reference convention differ from the corrected
v3 study, which uses \SIrange{30}{250}{MeV} model support and reports the
\SIrange{50}{250}{MeV} search range in
Figure~\ref{fig:v3-lslb-closure-2021}.}
\label{fig:closure-validation-2021-funcform}
\end{figure}

\FloatBarrier

\subsection{Earlier two-panel matched-reference displays}

The following pairs are the former dataset-level pull and residual-significance
summaries. They document the introduction of the matched-reference full-refit
procedure, but they were not generated as a common lower-bound comparison and are
therefore auxiliary to the v3 choice. In these studies ``matched reference'' means
that, at every mass, the injection was defined by
$A_{\mathrm{inj}}=Z_{\mathrm{inj}}\sigma_{A,\mathrm{ref}}$ using a background-only
refit with the same geometry, kernel policy, and extraction prescription as the
subsequent injected-toy fit. This removes a change of uncertainty convention between
signal definition and recovery.

\begin{figure}[p]
\centering
\graphicorplaceholder{0.84\linewidth}{injection_extraction_figs/refmatched_train2p25_pull_vs_mass_2015.png}

\vspace{0.35em}

\graphicorplaceholder{0.84\linewidth}{injection_extraction_figs/refmatched_train2p25_z_calibration_residual_2015.png}
\caption{Historical 2015 train-$2.25$ matched-reference functional-form closure
diagnostics. The injected amplitude was defined with the matched background-only
reference scale, after which every toy was processed with the same GP refit and
widened-blind prescription. In the upper panel, the zero-injection row tests for a
spurious fitted signal, the pull means test estimator bias, and the pull widths compare
the fitted $\sigma_A$ with the toy-to-toy scatter. The lower panel tests significance
calibration for the injected rows. This study established the matched-reference
workflow, while the explicit lower-bound scan in
Figure~\ref{fig:v3-lslb-closure-2015} determines the v3 2015 choice.}
\label{fig:refmatched-train225-closure-2015}
\end{figure}

\begin{figure}[p]
\centering
\graphicorplaceholder{0.84\linewidth}{injection_extraction_figs/refmatched_train2p25_pull_vs_mass_2016.png}

\vspace{0.35em}

\graphicorplaceholder{0.84\linewidth}{injection_extraction_figs/refmatched_train2p25_z_calibration_residual_2016.png}
\caption{Historical 2016 10\% train-$2.25$ matched-reference functional-form closure
diagnostics. The upper panel tests estimator closure and uncertainty calibration and
the lower panel tests significance calibration. The injections use the matched
background-only reference scale and each pseudoexperiment is subjected to the full GP
refit, so the plots probe both the local extraction and possible absorption by the
sideband-trained background. This single-bound study does not replace the dedicated
$k_{\min}=0.9$ comparison in Figure~\ref{fig:v3-lslb-closure-2016}.}
\label{fig:refmatched-train225-closure-2016}
\end{figure}

\begin{figure}[p]
\centering
\graphicorplaceholder{0.84\linewidth}{injection_extraction_figs/refmatched_train2p25_pull_vs_mass_2021.png}

\vspace{0.35em}

\graphicorplaceholder{0.84\linewidth}{injection_extraction_figs/refmatched_train2p25_z_calibration_residual_2021.png}
\caption{Historical 2021 1\% train-$2.25$ matched-reference functional-form closure
diagnostics. The layout and matched-reference convention mirror the 2015 and 2016
studies, allowing the fitted pull calibration and $\hat Z-Z_{\mathrm{inj}}$ response
to be compared across campaigns. This plot used the earlier 2021 input and support.
The corrected lower-bound study in Figure~\ref{fig:v3-lslb-closure-2021} instead uses
\SIrange{30}{250}{MeV} model support and reports the
\SIrange{50}{250}{MeV} search range.}
\label{fig:refmatched-train225-closure-2021}
\end{figure}

\FloatBarrier

\subsection{Legacy limit and projection figures}

The figures in this subsection are retained to make changes between review iterations
traceable. They do not use the complete July 8 v3 configuration and should not be
quoted as the present exclusion or sensitivity. The first five plots use the former
95\% CL convention; they are separated from the later 90\% CL projections so that the
two confidence levels cannot be mistaken for one common result series.

\begin{figure}[p]
\centering
\graphicorplaceholder{0.90\linewidth}{final_limit_projection_figs/individual_datasets/2015_eps2_coupling_bands_observed_expected.png}
\caption{Historical individual 2015 95\% CL upper limit on $\eps^2$ obtained with
\CLs{}. The observed curve is shown with the background-only expected median and its
68\% and 95\% intervals. This plot checked the 2015 yield-to-coupling conversion and
served as an input diagnostic for the earlier shared-coupling analysis; it is not a
component of the v3 90\% CL result.}
\label{fig:app-legacy-95-individuals}
\end{figure}

\begin{figure}[p]
\centering
\graphicorplaceholder{0.90\linewidth}{final_limit_projection_figs/individual_datasets/2016_eps2_coupling_bands_observed_expected.png}
\caption{Historical individual 2016 10\% 95\% CL upper limit on $\eps^2$ obtained
with \CLs{}. The observed curve and background-only expected bands were used to audit
the intermediate-mass input to an earlier combination. The analyzed sample is the
10\% development subset, not the complete 2016 campaign dataset, and this plot does
not use the v3 90\% CL configuration.}
\label{fig:app-legacy-95-individual-2016}
\end{figure}

\begin{figure}[p]
\centering
\graphicorplaceholder{0.90\linewidth}{final_limit_projection_figs/individual_datasets/2021_eps2_coupling_bands_observed_expected.png}
\caption{Historical individual 2021 1\% 95\% CL upper limit on $\eps^2$ obtained
with \CLs{}. The plot records the earlier 2021 input and conversion used during
development of the combined likelihood. It predates the regenerated
\texttt{final\_1pct\_invM.root} sample and is therefore retained only as analysis history,
not as a cross-check of the corrected v3 2021 limit.}
\label{fig:app-legacy-95-individual-2021}
\end{figure}

\begin{figure}[p]
\centering
\graphicorplaceholder{0.90\linewidth}{final_limit_projection_figs/combined/combined_corrected2_bands_with_correct_observed_eps2.png}
\caption{Historical simultaneous 2015$+$2016 10\%$+$2021 1\% 95\% CL upper limit
on $\eps^2$ obtained with \CLs{}. The green and yellow regions are the 68\% and 95\%
background-only expected intervals and the black curve is the then-current observed
shared-coupling limit. This plot documented the earlier combined likelihood; it has
been replaced by the selected v3 90\% CL configuration.}
\label{fig:app-legacy-95-combined}
\end{figure}

\begin{figure}[p]
\centering
\graphicorplaceholder{0.90\linewidth}{final_limit_projection_figs/combined/observed_eps2_overlay_individual_and_correct_combined.png}
\caption{Historical observed 95\% CL upper limits on $\eps^2$ for the individual
datasets and their simultaneous shared-coupling combination. The plot was used to
check that the joint likelihood behaved recognizably relative to its inputs; the
combined curve is a direct profile of one common $\eps^2$, not the pointwise envelope
of the individual curves. It is superseded by the clean v3 90\% CL comparison in the
main text.}
\label{fig:app-legacy-95-observed-overlay}
\end{figure}

The following two figures are 90\% CL projection products, but they predate the frozen
$k_{\min}=(1.0,0.9,1.1)$ configuration. They were produced by rescaling the current
2016 10\% and 2021 1\% density inputs to the complete campaign datasets while keeping
2015 fixed. They are observed-equivalent scaling studies, not likelihood fits to the
blinded data and not median expected sensitivities.

\begin{figure}[p]
\centering
\graphicorplaceholder{0.90\linewidth}{final_limit_projection_figs/90cls/projections/projected_full_unblinded_reach_90cls_eps2.png}
\caption{Historical 90\% CL projection to the complete 2016 and 2021 campaign
datasets. The 2016 10\% and 2021 1\% density inputs were scaled to their corresponding
complete samples while the fully analyzed 2015 input was held fixed. Because no
complete-dataset spectrum was fitted, the curve is an observed-equivalent projection
and should not be read as an expected sensitivity or observed exclusion. The v3
counterpart is Figure~\ref{fig:results-v3-projection}.}
\label{fig:app-legacy-projections}
\end{figure}

\begin{figure}[p]
\centering
\graphicorplaceholder{0.90\linewidth}{final_limit_projection_figs/90cls/babar_90_comparison/babar90_vs_projected_full_data_100pct.png}
\caption{Historical comparison of the projected HPS 90\% CL reach with the BaBar
90\% CL contour. Only the projection is shown on the HPS side; the observed HPS curve
was intentionally omitted. The lower panel gives the projected HPS limit divided by
the BaBar contour. The figure was used to locate the mass range in which a complete
HPS campaign could add sensitivity, not to quote a current exclusion.}
\label{fig:app-legacy-projection-babar}
\end{figure}

\FloatBarrier

\subsection{Plot-repair and signal-model diagnostics}

Plot-level repairs were useful for locating export inconsistencies but are not a
substitute for production reruns. Figures~\ref{fig:app-v2-repair-diagnostics} and
\ref{fig:app-v2-tail-diagnostic} record the repair procedure and limit-tail diagnostic
used for an earlier v2 export; their repaired mass points do not describe the July 8
candidate. The corresponding machine-readable audit tables for the v3 figures are
bundled with their provenance files and are summarized in
Section~\ref{sec:results-frozen-candidate}.

\begin{figure}[p]
\centering
\graphicorplaceholder{0.94\linewidth}{final_limit_projection_figs/90cls/diagnostics/combined_expected_band_spike_diagnostics_90cls.png}
\caption{Historical v2 expected-band repair diagnostic for the simultaneous 90\% CL
upper limit. The original expected quantiles are compared with the plotting values
after the isolated coherent excursions at 103 and 198~MeV were replaced by log-linear
interpolation between the adjacent mass hypotheses with the same active dataset set.
Only expected-band quantiles were changed; no observed limit or local-significance
value was modified. The plot documents how the display was repaired, rather than
providing statistical evidence for the interpolation.}
\label{fig:app-v2-repair-diagnostics}
\end{figure}

\begin{figure}[p]
\centering
\graphicorplaceholder{0.94\linewidth}{final_limit_projection_figs/90cls/diagnostics/combined_toy_limit_tail_areas_90cls.png}
\caption{Historical v2 limit-tail diagnostic for the simultaneous 90\% CL upper
limit. The connected curves give, at each mass, the fractions of background-only toy
limits at least as strong as, at least as weak as, and at least as extreme on either
side of the median as the observed limit. These are limit-ensemble tail probabilities,
not discovery $p_0$ values. The horizontal references show the corresponding local
and \v{S}id\'ak effective-trials thresholds used for review-level scan interpretation.}
\label{fig:app-v2-tail-diagnostic}
\end{figure}

The next three figures are historical individual discovery scans. In each case the
solid curve is the one-sided local $p_0(m)$ for a positive signal at a fixed mass; the
global curve applies the analytic effective-trials conversion used in that review
iteration. The separate panels are kept large enough to inspect the mass dependence
and are not interchangeable with the toy-limit tail probabilities above.

\begin{figure}[p]
\centering
\graphicorplaceholder{0.90\linewidth}{final_limit_projection_figs/individual_datasets/2015_local_p0_with_global_p_overlay.png}
\caption{Historical 2015 one-sided local discovery $p_0$ scan with the corresponding
effective-trials global approximation. The local minimum identifies the largest
upward fluctuation in the fully analyzed 2015 sample; the global curve accounts for
the mass scan under the approximation documented in Section~\ref{sec:methodology}.}
\label{fig:local-significances-across-datasets}
\end{figure}

\begin{figure}[p]
\centering
\graphicorplaceholder{0.90\linewidth}{final_limit_projection_figs/individual_datasets/2016_local_p0_with_global_p_overlay.png}
\caption{Historical 2016 10\% one-sided local discovery $p_0$ scan with the
corresponding effective-trials global approximation. The curve was used to locate
upward fluctuations in the development subset and to identify where the individual
spectrum could weaken the simultaneous observed limit.}
\label{fig:app-historical-local-p0-2016}
\end{figure}

\begin{figure}[p]
\centering
\graphicorplaceholder{0.90\linewidth}{final_limit_projection_figs/individual_datasets/2021_local_p0_with_global_p_overlay.png}
\caption{Historical 2021 1\% one-sided local discovery $p_0$ scan with the
corresponding effective-trials global approximation. This figure uses the earlier
2021 input and is retained to document development of the combined significance
workflow; the corrected individual scan in the main results section supersedes it.}
\label{fig:app-historical-local-p0-2021}
\end{figure}

\begin{figure}[p]
\centering
\graphicorplaceholder{0.90\linewidth}{final_limit_projection_figs/diagnostics/combined_local_p0_with_global_p_overlay.png}
\caption{Historical one-sided local $p_0$ scan for an earlier three-dataset
combination, together with the analytic effective-trials global approximation. The
local curve is the fixed-mass profile-likelihood discovery test; the global curve
converts its scan minimum to a review-level look-elsewhere estimate. The updated v3
combined curve is shown in Figure~\ref{fig:results-v3-local-p0}.}
\label{fig:combined-significance-rpen}
\end{figure}

\begin{figure}[p]
\centering
\graphicorplaceholder{0.90\linewidth}{final_limit_projection_figs/diagnostics/combined_local_Z_with_global_Z_overlay.png}
\caption{Gaussian-equivalent representation of the historical combined local and
effective-trials global significance curves. It contains the same hypothesis tests as
Figure~\ref{fig:combined-significance-rpen}, expressed as one-sided normal
significances rather than $p$-values. The pair was used to check the numerical
$p$-to-$Z$ conversion and to communicate the size of scan-level fluctuations.}
\label{fig:app-combined-significance-z}
\end{figure}

\begin{figure}[p]
\centering
\graphicorplaceholder{0.96\linewidth}{final_limit_projection_figs/90cls/scaled_21/scaled_21_bandfixed_dimuon_vs_unscaled_overlay_ratio.png}
\caption{Historical combined 90\% CL upper-limit comparison used to study the 2021
mass-resolution correction. Blue curves use the unscaled \texttt{v2\_dimuon}
reference and red curves use the target-constrained 2021 resolution scale under the
same dimuon-threshold branching convention. The lower panel gives scaled/unscaled
ratios for the observed limit and expected median; at common masses the branching
factor cancels, isolating changes caused by the broadened 2021 signal template and
shared-coupling fit. The effect was not monotonic in mass, so this plot supports a
signal-model cross-check rather than a uniform correction factor. It is not the
corrected 2021 result shown in Figure~\ref{fig:results-v3-2021-pair}.}
\label{fig:app-2021-resolution-comparison}
\end{figure}

\FloatBarrier

\subsection{Auxiliary shared-coupling demonstrations}

The following plots illustrate shared-coupling bookkeeping but were generated with
provisional injection scenarios. They are not direct \CLs{} coverage tests and are not
used to select the v3 kernel bounds. Their purpose is narrower: they show how one
injected $\eps^2$ maps through the dataset-specific conversion factors into different
signal yields and significances, and how statistically independent likelihoods add
information about that common parameter.

\begin{figure}[p]
\centering
\graphicorplaceholder{0.90\linewidth}{combined_search_figs/combined_search_power_scenarios.png}
\caption{Historical two-dataset shared-coupling search-power scan. At each point the
injected common $\eps^2$ is translated into dataset-specific signal yields before the
combined statistic is evaluated. The curves therefore illustrate genuine accumulation
of likelihood information for one coupling, rather than a toy in which the same event
count is added to every dataset. The injection scenarios were provisional, so the
plot documents the combination logic but does not quantify the v3 candidate
sensitivity.}
\label{fig:app-combined-search-power}
\end{figure}

\begin{figure}[p]
\centering
\graphicorplaceholder{0.90\linewidth}{combined_search_figs/combined_search_power_constituent_pvalues_5sigma.png}
\caption{Historical constituent-dataset allocations capable of reaching a target
combined $5\sigma$ excess in the provisional shared-coupling study. The panel shows
that a fixed combined threshold can arise from different per-dataset contributions,
depending on the local resolution, background, and signal-normalization leverage.
It should be read with Figure~\ref{fig:app-combined-search-power} as an illustration
of the likelihood bookkeeping, not as a prediction of discovery reach.}
\label{fig:app-combined-search-power-constituents}
\end{figure}

The extraction displays below use no-refit pseudoexperiments with a nominal
five-standard-deviation shared signal. The GP background is first conditioned on the
sidebands and then held fixed while the local pseudoexperiment is extracted. This
isolates the signal and nuisance-parameter bookkeeping of the local likelihood; it
does not test whether a toy-by-toy sideband refit can absorb signal.

The final display in this sequence answers that complementary question. Its toy was
generated over the full mass range and the GP was retrained on the toy sidebands before
the blind-window extraction. Any loss of recovery can therefore reveal signal
absorption by the background refit, which the three no-refit displays are not designed
to measure.

\clearpage
\begin{landscape}
\begin{figure}[H]
\centering
\graphicorplaceholder{0.92\linewidth}{extract_display_combined_m040MeV_z5p0_polished.png}
\caption{Historical no-refit shared-coupling extraction at \SI{40}{MeV}. The upper
panels show the dataset-level pseudoexperiments, conditioned GP backgrounds, and
extracted signal components. The lower panel places those components on the common
$(m-m_0)/\sigma_m$ coordinate used for the simultaneous interpretation. The annotation
records the injected and fitted coupling, combined upper limit, local significance,
and per-dataset fit summaries. In this low-mass example the 2015 sample carries most
of the signal normalization, while the fit still constrains one common $\eps^2$.}
\label{fig:app-combined-extract-no-refit}
\end{figure}
\end{landscape}

\begin{landscape}
\begin{figure}[H]
\centering
\graphicorplaceholder{0.92\linewidth}{extract_display_combined_m080MeV_z5p0_polished.png}
\caption{Historical no-refit shared-coupling extraction at \SI{80}{MeV}. The same
five-standard-deviation injection and held-fixed sideband-conditioned background
prescription are used as in Figure~\ref{fig:app-combined-extract-no-refit}. The
relative dataset contributions differ from the \SI{40}{MeV} example because the
accepted statistics, mass resolution, and prompt-background level vary with campaign
and mass.}
\label{fig:app-combined-extract-no-refit-80}
\end{figure}
\end{landscape}

\begin{landscape}
\begin{figure}[H]
\centering
\graphicorplaceholder{0.92\linewidth}{extract_display_combined_m120MeV_z5p0_polished.png}
\caption{Historical no-refit shared-coupling extraction at \SI{120}{MeV}. The display
uses the same provisional five-standard-deviation shared signal as the lower-mass
examples. At this mass the changing active samples and conversion factors redistribute
the fitted signal among the datasets, illustrating why the simultaneous likelihood
must map a common coupling to campaign-dependent yields instead of adding separate
limits after fitting.}
\label{fig:app-combined-extract-no-refit-120}
\end{figure}
\end{landscape}

\begin{landscape}
\begin{figure}[H]
\centering
\graphicorplaceholder{0.92\linewidth}{extract_display_combined_m040MeV_z3p0_refit.png}
\caption{Historical full-refit shared-coupling extraction at \SI{40}{MeV} for a
provisional three-standard-deviation injection. The full pseudoexperiment was
regenerated and the GP background was retrained on its sidebands before extraction.
The visibly poorer recovery documented signal absorption in that provisional
configuration. It is retained as a failure-mode diagnostic and is not evidence of
calibrated v3 signal recovery.}
\label{fig:app-combined-extract-refit}
\end{figure}
\end{landscape}
\clearpage
